\documentclass[12pt,a4paper]{article}

%!TEX program = pdflatex

% =============================================================================
% 1. PHÔNG CHỮ & NGÔN NGỮ (Hỗ trợ Tiếng Việt)
% =============================================================================
\usepackage[utf8]{inputenc}
\usepackage[T5]{fontenc}
\usepackage[vietnamese]{babel}

% =============================================================================
% 2. BỐ CỤC TRANG (LAYOUT)
% =============================================================================
\usepackage[margin=2.5cm]{geometry}
\usepackage{setspace}
\onehalfspacing

% =============================================================================
% 3. ĐỒ HỌA & BIỂU ĐỒ
% =============================================================================
\usepackage{graphicx}
\usepackage{float}

% =============================================================================
% 4. BẢNG BIỂU & DANH SÁCH
% =============================================================================
\usepackage{array}
\usepackage{booktabs}
\usepackage{longtable}
\usepackage{enumitem}
\usepackage{tabularx}

% =============================================================================
% 5. TIỆN ÍCH VĂN BẢN & LIÊN KẾT
% =============================================================================
\usepackage{xcolor}
\usepackage{titlesec}
\usepackage{tocloft}
\usepackage{xurl}
\usepackage{seqsplit}
\usepackage{hyperref}

\hypersetup{
    colorlinks=false,
    hidelinks,
    pdftitle={Tổng quan và thu thập yêu cầu hệ thống TMS},
    pdfauthor={Nhóm thực hiện}
}

% =============================================================================
% 6. ĐỊNH DẠNG ĐOẠN VĂN & ĐÁNH SỐ MỤC
% =============================================================================
\setlength{\parindent}{0pt}
\setlength{\parskip}{6pt}
\setcounter{secnumdepth}{3}

\renewcommand{\thesection}{\Roman{section}}
\renewcommand{\thesubsection}{\arabic{subsection}}
\renewcommand{\thesubsubsection}{\thesubsection.\arabic{subsubsection}}

\titleformat{\section}{\large\bfseries}{\thesection.}{0.5em}{}
\titleformat{\subsection}{\normalsize\bfseries}{\thesubsection.}{0.5em}{}
\titleformat{\subsubsection}{\normalsize\itshape\bfseries}{\thesubsubsection.}{0.5em}{}

\renewcommand{\cftsecaftersnum}{.}
\renewcommand{\cftsubsecaftersnum}{.}
\renewcommand{\cftsubsubsecaftersnum}{.}

% =============================================================================
% 7. CUSTOM COMMANDS & COLUMN TYPES
% =============================================================================
\newcolumntype{C}[1]{>{\centering\arraybackslash}p{#1}}
\newcolumntype{L}[1]{>{\raggedright\arraybackslash}p{#1}}
\newcolumntype{Y}{>{\raggedright\arraybackslash}X}
\newcommand{\code}[1]{{\ttfamily\seqsplit{#1}}}
\newcommand{\tablesetup}{\small\setlength{\tabcolsep}{3pt}\renewcommand{\arraystretch}{1.2}}

\begin{document}

\section{Tổng quan đề tài}

\subsection{Giới thiệu đề tài}

Trong hoạt động dạy học theo mô hình lớp nhỏ hoặc gia sư, giáo viên thường phải
quản lý đồng thời nhiều nhóm thông tin: danh sách học sinh, lớp học, lịch học,
buổi học thực tế, điểm danh, học phí, giao dịch tài chính và kết quả luyện tập.
Nếu các dữ liệu này được ghi nhận rời rạc bằng sổ tay, bảng tính hoặc tin nhắn
riêng lẻ, việc theo dõi tình hình học tập và vận hành lớp học dễ phát sinh sai
sót, trùng lặp và thiếu tính kịp thời.

Bên cạnh đó, nhiều hoạt động hỗ trợ học tập hiện nay diễn ra trên các nền tảng
ngoài hệ thống chính. Discord thường được dùng để tổ chức kênh trao đổi, gửi
thông báo và hỗ trợ điểm danh qua kênh thoại. Codeforces Gym được dùng để
giao bài luyện tập và theo dõi bảng xếp hạng. Khi các dữ liệu này không được
kết nối với hệ thống quản lý lớp học, giáo viên phải đối chiếu thủ công giữa
nhiều nguồn khác nhau, làm tăng thời gian vận hành và giảm khả năng nắm bắt
tình hình thực tế.

Vì vậy, đề tài ``Tutor Management System'' (TMS) được thực hiện nhằm xây dựng
một hệ thống web tập trung cho nghiệp vụ quản lý dạy học. Hệ thống hỗ trợ giáo
viên quản lý học sinh, lớp học, buổi học, điểm danh, tài chính, đồng thời tích
hợp Discord và Codeforces để giảm thao tác thủ công trong các quy trình thường
xuyên. Đối với quản trị viên, hệ thống cung cấp các chức năng quản lý tài khoản
giáo viên và cấu hình bot Discord dùng chung.

\subsection{Mục tiêu phần mềm}

Phần mềm được xây dựng nhằm đạt được các mục tiêu cốt lõi sau:

\begin{itemize}[leftmargin=*]
    \item \textbf{Tin học hóa nghiệp vụ quản lý dạy học:} Tập trung hóa dữ liệu
    về giáo viên, học sinh, lớp học, lịch học, buổi học và quá trình ghi danh
    trên một hệ thống thống nhất.

    \item \textbf{Theo dõi học sinh và lớp học chính xác:} Cho phép giáo viên
    quản lý hồ sơ học sinh, trạng thái ghi danh, lịch sử học tập, lớp đang học
    và các thay đổi như chuyển lớp, rút khỏi lớp hoặc lưu trữ học sinh.

    \item \textbf{Hỗ trợ điểm danh và quản lý buổi học:} Cung cấp chức năng tạo
    buổi học, cập nhật trạng thái buổi học và ghi nhận điểm danh để giáo viên
    có dữ liệu phục vụ theo dõi học tập và tính phí.

    \item \textbf{Quản lý tài chính học sinh:} Ghi nhận giao dịch, khoản phí
    phát sinh từ buổi học, số dư của học sinh và các báo cáo tài chính phục vụ
    đối soát.

    \item \textbf{Tích hợp công cụ học tập và giao tiếp:} Kết nối Discord để hỗ
    trợ thiết lập máy chủ/kênh, gửi thông báo và liên kết định danh; kết nối
    Codeforces để quản lý Gym, bài tập và bảng xếp hạng luyện tập.

    \item \textbf{Hỗ trợ ra quyết định:} Cung cấp màn hình tổng quan, báo cáo
    học sinh, báo cáo tài chính và dữ liệu kết quả học tập để giáo viên nhanh
    chóng nắm bắt tình hình lớp học.

    \item \textbf{Kiểm soát truy cập và bảo mật:} Thiết lập cơ chế đăng nhập,
    phân quyền theo vai trò quản trị viên và giáo viên, đồng thời giới hạn giáo
    viên chỉ truy cập dữ liệu thuộc phạm vi quản lý của mình.
\end{itemize}

\subsection{Phạm vi phần mềm}

Phạm vi của phần mềm tập trung vào các chức năng cần thiết cho một hệ thống web
quản lý hoạt động dạy học. Hệ thống được phát triển theo các nhóm chức năng
chính sau:

\begin{itemize}[leftmargin=*]
    \item \textbf{Quản lý tài khoản và phân quyền:} Hỗ trợ đăng nhập, quản lý
    phiên làm việc, cập nhật hồ sơ cá nhân, quản lý tài khoản giáo viên và phân
    quyền giữa quản trị viên với giáo viên.

    \item \textbf{Quản lý học sinh:} Cho phép thêm, cập nhật, tra cứu, lưu trữ,
    khôi phục học sinh; quản lý thông tin liên hệ, định danh Codeforces, trạng
    thái học tập và liên kết Discord của học sinh.

    \item \textbf{Quản lý ghi danh:} Theo dõi quan hệ giữa học sinh và lớp học,
    hỗ trợ chuyển lớp, rút khỏi lớp và lưu lịch sử ghi danh theo thời gian.

    \item \textbf{Quản lý lớp học và lịch học:} Cho phép tạo lớp, cập nhật thông
    tin lớp, quản lý lịch học cố định, tạo buổi học từ lịch và theo dõi trạng
    thái từng buổi học.

    \item \textbf{Quản lý điểm danh:} Ghi nhận kết quả điểm danh của học sinh
    trong từng buổi học, phục vụ theo dõi tham gia học tập và phát sinh khoản
    phí liên quan.

    \item \textbf{Quản lý tài chính:} Quản lý giao dịch, khoản phí, số dư học
    sinh và báo cáo tài chính theo phạm vi dữ liệu của giáo viên.

    \item \textbf{Tích hợp Discord:} Quản lý cấu hình bot Discord, liên kết định
    danh giáo viên và học sinh, thiết lập Discord cho lớp học, chọn máy chủ/kênh
    Discord và gửi thông báo đến học sinh.

    \item \textbf{Tích hợp Codeforces Gym:} Quản lý thông tin Gym, gắn Gym cho
    lớp học, đồng bộ bài tập và bảng xếp hạng để giáo viên theo dõi kết quả
    luyện tập.

    \item \textbf{Báo cáo và màn hình tổng quan:} Cung cấp màn hình tổng quan, báo cáo
    học sinh, báo cáo tài chính và bảng xếp hạng Gym giúp giáo viên quan sát
    nhanh tình hình vận hành.
\end{itemize}

Trong phạm vi hiện tại, hệ thống tập trung vào ứng dụng web cho giáo viên và
quản trị viên. Hệ thống chưa bao gồm ứng dụng di động riêng, cổng thanh toán tự
động, lớp học trực tuyến/video meeting, hệ thống chấm bài thay thế Codeforces
hoặc cổng tự phục vụ đầy đủ cho học sinh. Các nội dung này có thể được xem là
hướng mở rộng trong các giai đoạn phát triển tiếp theo.

\section{Xác định và thu thập yêu cầu}

\subsection{Kế hoạch và phương pháp khảo sát}

\subsubsection{Đối tượng và địa điểm khảo sát}

\begin{itemize}[leftmargin=*]
    \item \textbf{Bối cảnh khảo sát:} Khảo sát được định hướng trên mô hình lớp
    học/gia sư quy mô nhỏ, trong đó giáo viên quản lý nhiều lớp, nhiều học sinh,
    nhiều buổi học và nhiều khoản thu chi phát sinh. Hoạt động hiện tại thường
    kết hợp bảng tính, tin nhắn riêng, Discord và Codeforces nên dữ liệu bị phân
    tán ở nhiều nơi.

    \item \textbf{Giáo viên:} Là người dùng chính của hệ thống, trực tiếp quản lý
    học sinh, lớp học, lịch học, buổi học, điểm danh, giao dịch tài chính, Discord
    của lớp và Codeforces Gym.

    \item \textbf{Quản trị viên:} Là người quản lý tài khoản giáo viên, kiểm soát
    trạng thái tài khoản và cấu hình bot Discord dùng chung cho hệ thống.

    \item \textbf{Học sinh:} Là đối tượng được giáo viên quản lý. Học sinh tham
    gia lớp học, có lịch sử ghi danh, kết quả điểm danh, số dư tài chính, thông
    tin Discord và kết quả luyện tập Codeforces.

    \item \textbf{Hệ thống ngoài:} Discord và Codeforces là các nguồn dữ liệu,
    kênh giao tiếp và công cụ học tập cần được khảo sát để xác định phạm vi tích
    hợp phù hợp.
\end{itemize}

\subsubsection{Các phương pháp khảo sát}

\begin{itemize}[leftmargin=*]
    \item \textbf{Phỏng vấn:} Thu thập thông tin từ giáo viên và quản trị viên về
    các khó khăn khi quản lý học sinh, lớp học, lịch học, điểm danh, học phí,
    Discord và Codeforces bằng nhiều công cụ rời rạc.

    \item \textbf{Nghiên cứu tài liệu nghiệp vụ:} Phân tích các dạng dữ liệu mà
    hệ thống cần thay thế hoặc tập trung hóa, gồm danh sách học sinh, danh sách
    lớp học, lịch học, bảng điểm danh, bảng giao dịch, bảng công nợ, cấu hình
    Discord và bảng xếp hạng Codeforces Gym.

    \item \textbf{Quan sát quy trình hiện tại:} Xem xét các thao tác thường gặp
    như tạo lớp, ghi danh học sinh, tạo buổi học, điểm danh, ghi nhận khoản phí,
    gửi thông báo lớp học và kiểm tra kết quả luyện tập trên Codeforces.

    \item \textbf{Bảng câu hỏi:} Dùng để xác nhận mức độ ưu tiên của các chức
    năng như màn hình tổng quan, quản lý học sinh, quản lý lớp, điểm danh, báo cáo tài
    chính, tích hợp Discord và theo dõi Codeforces Gym.
\end{itemize}

\subsection{Khảo sát hiện trạng}

\subsubsection{Hiện trạng tổ chức}

Cơ cấu vận hành của hệ thống mục tiêu gồm các nhóm vai trò chính sau:

\begin{itemize}[leftmargin=*]
    \item \textbf{Quản trị viên:} Quản lý tài khoản giáo viên, kích hoạt hoặc vô
    hiệu hóa tài khoản, thiết lập thông tin bot Discord và theo dõi cấu hình hệ
    thống dùng chung.

    \item \textbf{Giáo viên:} Quản lý toàn bộ nghiệp vụ dạy học trong phạm vi của
    mình. Mỗi giáo viên có thể quản lý nhiều học sinh, nhiều lớp, nhiều buổi học,
    nhiều giao dịch và nhiều Codeforces Gym.

    \item \textbf{Học sinh:} Được ghi danh vào lớp, tham gia các buổi học, phát
    sinh điểm danh và khoản phí, đồng thời có thể liên kết Discord hoặc Codeforces
    để phục vụ giao tiếp và luyện tập.
\end{itemize}

Về phân quyền, hệ thống cần tách rõ dữ liệu quản trị và dữ liệu nghiệp vụ. Quản
trị viên chỉ xử lý tài khoản giáo viên và cấu hình hệ thống; giáo viên chỉ được
truy cập học sinh, lớp học, buổi học, giao dịch, cấu hình liên kết Discord và Gym thuộc
phạm vi sở hữu của mình.

\subsubsection{Hiện trạng nghiệp vụ}

Kết quả khảo sát cho thấy các quy trình nghiệp vụ cốt lõi có quan hệ chặt chẽ
với nhau nhưng thường bị xử lý trên nhiều công cụ khác nhau. Điều này làm giáo
viên mất thời gian đối chiếu và khó nhìn thấy bức tranh tổng thể của từng học
sinh hoặc từng lớp học.

\paragraph{2.2.1. Quy trình quản lý học sinh và ghi danh}

Quy trình này thường diễn ra theo trình tự sau:

\begin{itemize}[leftmargin=*]
    \item Giáo viên ghi nhận thông tin học sinh như họ tên, thông tin liên hệ,
    trạng thái học tập và định danh Codeforces.
    \item Khi học sinh tham gia lớp, giáo viên ghi nhận thông tin ghi danh vào
    lớp tương ứng.
    \item Nếu học sinh chuyển lớp hoặc ngừng học, giáo viên phải cập nhật lại
    trạng thái ghi danh và lịch sử học tập.
    \item Khi học sinh tạm ngừng hoặc không còn theo học, giáo viên cần lưu trữ
    hồ sơ để dữ liệu cũ vẫn được bảo toàn cho báo cáo và đối soát.
\end{itemize}

Nếu quy trình này được thực hiện bằng bảng tính riêng lẻ, giáo viên dễ gặp tình
trạng một học sinh xuất hiện ở nhiều nơi, trạng thái ghi danh không đồng nhất và
khó truy xuất lịch sử thay đổi.

\paragraph{2.2.2. Quy trình quản lý lớp học, lịch học và buổi học}

\begin{itemize}[leftmargin=*]
    \item Giáo viên tạo lớp học, cấu hình học phí mỗi buổi và thiết lập lịch học
    cố định trong tuần.
    \item Từ lịch học, giáo viên cần tạo hoặc theo dõi các buổi học cụ thể theo
    ngày giờ thực tế.
    \item Khi có thay đổi, buổi học có thể được cập nhật trạng thái hoặc hủy.
    \item Mỗi buổi học cần liên kết với danh sách học sinh được ghi danh để phục
    vụ điểm danh và phát sinh phí.
\end{itemize}

Trong cách làm thủ công, lịch học, buổi học và danh sách học sinh thường bị tách
thành nhiều bảng hoặc nhiều tin nhắn. Khi lớp thay đổi lịch hoặc phát sinh buổi
học đặc biệt, việc đồng bộ lại dữ liệu mất nhiều thời gian.

\paragraph{2.2.3. Quy trình điểm danh và tài chính}

\begin{itemize}[leftmargin=*]
    \item Giáo viên mở buổi học và ghi nhận tình trạng tham gia của từng học
    sinh.
    \item Kết quả điểm danh được dùng làm căn cứ phát sinh khoản phí hoặc điều
    chỉnh khoản phí cho học sinh.
    \item Giáo viên ghi nhận các giao dịch tiền vào/tiền ra của học sinh để theo
    dõi số dư.
    \item Cuối kỳ hoặc khi cần đối soát, giáo viên tổng hợp khoản phí, giao dịch
    và số dư để lập báo cáo tài chính.
\end{itemize}

Nếu điểm danh và tài chính không nằm trong cùng một hệ thống, giáo viên phải
đối chiếu thủ công giữa buổi học, trạng thái tham gia và giao dịch. Việc này dễ
dẫn đến sai lệch số dư, ghi thiếu phí hoặc khó giải thích lịch sử thay đổi cho
từng học sinh.

\paragraph{2.2.4. Quy trình tích hợp Discord và Codeforces}

\begin{itemize}[leftmargin=*]
    \item Với Discord, giáo viên cần liên kết định danh, chọn máy chủ/kênh cho
    lớp học, gửi thông báo và quản lý thông tin học sinh liên quan Discord.
    \item Với Codeforces, giáo viên cần gắn Gym cho lớp, đồng bộ bài tập và theo
    dõi bảng xếp hạng của học sinh trong Gym.
    \item Các dữ liệu từ Discord và Codeforces cần được liên hệ lại với lớp học,
    học sinh và giáo viên tương ứng trong hệ thống TMS.
\end{itemize}

Khi thiếu tích hợp, giáo viên phải mở riêng Discord để gửi thông báo và mở riêng
Codeforces để xem kết quả luyện tập. Dữ liệu học sinh trong các nền tảng này
không tự động khớp với hồ sơ học sinh trong hệ thống quản lý lớp học.

\subsubsection{Hiện trạng tin học}

\textbf{2.3.1. Phần cứng đang sử dụng}

\begin{itemize}[leftmargin=*]
    \item Máy tính cá nhân hoặc laptop của giáo viên để quản lý bảng tính, truy
    cập Discord, Codeforces và hệ thống web.
    \item Điện thoại thông minh để trao đổi nhanh với học sinh, nhận thông báo
    và kiểm tra thông tin khi không ngồi trước máy tính.
    \item Kết nối Internet ổn định để truy cập hệ thống web và các nền tảng bên
    ngoài.
\end{itemize}

\textbf{2.3.2. Phần mềm đang sử dụng}

\begin{itemize}[leftmargin=*]
    \item Bảng tính như Excel hoặc Google Sheets để lưu danh sách học sinh, lịch
    học, điểm danh, giao dịch và công nợ.
    \item Discord để trao đổi với học sinh, tổ chức kênh lớp học và gửi thông
    báo.
    \item Codeforces để giao Gym, xem bài tập và theo dõi bảng xếp hạng luyện
    tập.
    \item Các ứng dụng nhắn tin khác để trao đổi nhanh, nhưng dữ liệu từ các ứng
    dụng này không được đồng bộ vào hệ thống quản lý.
\end{itemize}

\textbf{2.3.3. Trình độ công nghệ thông tin của người dùng}

\begin{itemize}[leftmargin=*]
    \item Giáo viên có thể sử dụng trình duyệt web, bảng tính, Discord và
    Codeforces ở mức phục vụ công việc hằng ngày.
    \item Quản trị viên có khả năng thao tác với tài khoản, cấu hình hệ thống và
    kiểm tra trạng thái hoạt động.
    \item Học sinh quen với việc sử dụng Discord và Codeforces, nhưng không phải
    là nhóm người dùng chính thao tác sâu trong hệ thống quản lý.
\end{itemize}

\subsection{Đánh giá hiện trạng và những tính năng mong muốn}

Từ kết quả khảo sát, nhóm tiến hành phân tích các ưu điểm, hạn chế và nhu cầu
chính để xác định yêu cầu cho hệ thống TMS.

\subsubsection{Ưu điểm}

\begin{itemize}[leftmargin=*]
    \item \textbf{Chi phí ban đầu thấp:} Bảng tính và các công cụ giao tiếp sẵn
    có không yêu cầu triển khai hệ thống phức tạp ngay từ đầu.

    \item \textbf{Linh hoạt:} Giáo viên có thể nhanh chóng thêm cột, sửa dòng,
    ghi chú hoặc điều chỉnh thông tin theo tình huống thực tế.

    \item \textbf{Dễ tiếp cận:} Người dùng đã quen với trình duyệt, bảng tính,
    Discord và Codeforces nên không cần thay đổi hoàn toàn thói quen làm việc.
\end{itemize}

\subsubsection{Nhược điểm và khó khăn cần giải quyết}

Đây là những vấn đề cốt lõi mà phần mềm mới cần giải quyết:

\textbf{a) Dữ liệu phân tán và thiếu đồng bộ}

Thông tin học sinh, lớp học, lịch học, điểm danh, giao dịch, Discord và
Codeforces nằm ở nhiều nơi khác nhau. Hậu quả là giáo viên khó nhìn thấy trạng
thái đầy đủ của một học sinh hoặc một lớp học tại một thời điểm.

\textbf{b) Sai sót trong điểm danh và tài chính}

Khi điểm danh và phát sinh phí được xử lý thủ công, hệ thống dễ ghi thiếu buổi,
ghi nhầm học sinh hoặc tính sai số dư. Việc truy lại nguyên nhân sai lệch cũng
mất nhiều thời gian vì dữ liệu không có lịch sử rõ ràng.

\textbf{c) Khó kiểm soát quyền truy cập}

Nếu nhiều giáo viên cùng sử dụng bảng tính hoặc thư mục chung, việc phân quyền
theo phạm vi sở hữu dữ liệu không rõ ràng. Dữ liệu học sinh và tài chính cần
được giới hạn theo giáo viên để bảo đảm bảo mật.

\textbf{d) Thiếu báo cáo tổng hợp}

Giáo viên cần biết nhanh số lượng học sinh, lớp đang hoạt động, buổi học gần
đây, khoản phí, giao dịch và tình hình luyện tập. Nếu dữ liệu nằm rời rạc, báo
cáo phải làm thủ công và thường có độ trễ.

\textbf{e) Tích hợp ngoài hệ thống còn thủ công}

Discord và Codeforces hỗ trợ trực tiếp cho hoạt động học tập nhưng thường được
theo dõi riêng. Giáo viên phải tự đối chiếu định danh học sinh, thông báo lớp và
kết quả Gym, dẫn đến tốn thời gian và dễ nhầm lẫn.

\subsubsection{Giải pháp và tính năng mong muốn}

Từ các khó khăn trên, nhóm tổng hợp được các tính năng người dùng mong muốn
nhất:

\begin{longtable}{L{3.2cm} L{11.2cm}}
\toprule
\textbf{Chủ thể} & \textbf{Mong muốn} \\
\midrule
\endhead
Giáo viên & ``Tôi muốn quản lý học sinh, lớp học, lịch học, buổi học, điểm danh và học phí trong cùng một hệ thống, không phải đối chiếu nhiều bảng tính.'' \\
Giáo viên & ``Tôi muốn khi điểm danh một buổi học, hệ thống có thể lưu kết quả và hỗ trợ theo dõi khoản phí, số dư của học sinh.'' \\
Giáo viên & ``Tôi muốn gắn Discord và Codeforces Gym vào lớp học để gửi thông báo và xem kết quả luyện tập nhanh hơn.'' \\
Quản trị viên & ``Tôi muốn quản lý tài khoản giáo viên, trạng thái hoạt động và cấu hình bot Discord ở một nơi rõ ràng.'' \\
Học sinh & ``Tôi muốn nhận thông báo lớp học qua Discord và được ghi nhận kết quả luyện tập Codeforces đúng với lớp đang học.'' \\
\bottomrule
\end{longtable}

\subsection{Xác định và thu thập yêu cầu}

Kế hoạch thu thập yêu cầu được thực hiện thông qua việc phỏng vấn các bên liên
quan chính, nghiên cứu dữ liệu nghiệp vụ hiện có và quan sát các quy trình
thường xuyên như quản lý học sinh, ghi danh, tạo buổi học, điểm danh, ghi nhận
giao dịch, gửi thông báo và theo dõi Codeforces Gym.

Các yêu cầu thu thập được được phân nhóm thành các mảng chức năng chính:

\begin{itemize}[leftmargin=*]
    \item Quản lý tài khoản, đăng nhập, phân quyền và phạm vi truy cập dữ liệu.
    \item Quản lý học sinh, hồ sơ học sinh, lưu trữ và khôi phục học sinh.
    \item Quản lý ghi danh, chuyển lớp, rút khỏi lớp và lịch sử học tập.
    \item Quản lý lớp học, lịch học cố định, buổi học và trạng thái buổi học.
    \item Quản lý điểm danh và dữ liệu tham gia buổi học.
    \item Quản lý giao dịch, khoản phí, số dư và báo cáo tài chính.
    \item Tích hợp Discord cho định danh, cấu hình lớp học và gửi thông báo.
    \item Tích hợp Codeforces Gym cho bài tập và bảng xếp hạng luyện tập.
    \item Cung cấp màn hình tổng quan và báo cáo giúp giáo viên theo dõi hoạt động dạy
    học một cách tập trung.
\end{itemize}

\end{document}
